\documentclass{article}
\usepackage{iclr2027_conference,times,graphicx,booktabs,amsmath}
\usepackage{colortbl,float}
\usepackage[colorlinks=true,linkcolor=black,citecolor=black]{hyperref}
\iclrfinalcopy
\newcommand{\ours}{\textsc{Pixel-OPSD}}
\title{Pixel-OPSD: Beyond Scalar Rewards for Pixel-Level MLLMs via On-Policy Self-Distillation}\author{Anonymous Authors}\begin{document}\setlength{\extrarowheight}{2pt}\raggedbottom\begingroup\hbadness=10000\maketitle\endgroup
\begin{abstract}CycleGRPO uses a scalar cycle reward for pixel-level multimodal generation, broadcasting the success or failure of a complete response across semantic and mask tokens. We present \ours{}, a training-time credit-assignment framework that progressively adds on-policy distillation (OPD), one spatial-credit module EGCA, and supervised anchors while preserving the actor, tokenizer, decoder, and inference interface. OPD densifies feedback over the prefixes actually sampled by the actor; EGCA aligns code-level updates with evidence from decoded partial masks; and supervised mixing stabilizes compositional language behavior that a narrow self-supervised cycle may forget. We evaluate the progression on GRES and GroundingSuite with a shared decoder, matched complete-data protocols, independent diagnostics, a progressive ablation, and coordinate-aligned qualitative cases. The resulting analysis separates semantic selection, spatial refinement, and language anchoring instead of reducing them to one terminal score.\end{abstract}
\section{Introduction}Pixel-level MLLMs must make a semantic decision and expose the pixels that justify it. A trajectory scalar cannot distinguish a wrong noun, relation, coarse code, or fine boundary, a limitation shared by scalar policy optimization such as PPO~\citep{shulman2017ppo}. \ours{} addresses this bottleneck progressively with OPD, EGCA, and supervised mixing while keeping the actor, tokenizer, and decoder fixed.

Pixel-level generation couples a semantic decision to a spatial program. In a coarse-to-fine code sequence, an early instance choice determines the region available to every later refinement. Yet the usual policy-gradient update broadcasts one terminal scalar over all tokens. A response may name the right object but miss its extent, or produce a plausible mask for the wrong instance; these errors have different causes but look identical to the optimizer.

This paper makes three contributions in strict progression. First, we adapt on-policy distillation (OPD) to pixel trajectories, turning the cycle outcome into dense token supervision while preserving sampled prefixes. Second, we introduce EGCA, the only new spatial-credit module, which uses decoded evidence to attribute overlap changes to mask codes. Third, we mix supervised examples with the self-supervised cycle to anchor semantic and refusal behavior. The contributions are separable: the first changes supervision granularity, the second spatial attribution, and the third the data stream.

We ask whether token-level OPD improves prefix consistency without collapsing exploration; whether EGCA localizes boundary and extent errors beyond a scalar reward; and whether supervised mixing stabilizes semantic behavior across benchmarks. The experiments, ablation, and qualitative panels answer these questions separately.
The separation is more than an implementation detail. Pixel trajectories contain at least three coupled decisions: selecting the referred instance, refining its extent, and expressing the answer in language that remains faithful to the query.

A single reward can improve one decision by sacrificing another, and an aggregate IoU cannot reveal that exchange. We therefore keep the actor's sampled response as the common object of study and change one source of supervision at a time.

This gives a direct correspondence between a failure mode and an intervention: OPD changes granularity, EGCA changes spatial attribution, and supervised mixing changes the data balance. The resulting protocol makes both positive and negative transfers informative, rather than hiding them behind a single rank.
\begin{figure*}[t]\centering\includegraphics[width=.98\textwidth]{figures/overview.pdf}\label{fig:overview}\caption{The overview separates the scalar-credit failure from our intervention. A trajectory-level cycle score aliases noun, relation, and boundary errors; OPD supplies token-level teaching, while EGCA localizes spatial credit to decoded mask codes.}\end{figure*}
\section{Related Work}
Recent pixel MLLMs combine autoregressive language with discrete mask tokens, enabling a single model to answer a referring question and emit a region~\citep{li2024pixellm,rasheed2024glamm,lai2024lisa,you2024ferret,yuan2024osprey}. SAMTok introduces a compact coarse-to-fine code interface~\citep{zhou2026samtok}, while concurrent systems improve mask decoding and grounding benchmarks~\citep{li2024omgseg,li2025rmpsam,wang2026gar,tian2026mmada}. These models expose a useful intermediate representation, but their training objectives generally supervise the final mask or language sequence rather than the individual decisions that produced it.

CycleGRPO supplies an on-policy cycle in which sampled text and mask codes are decoded, scored, and reinforced~\citep{zhang2026cyclegrpo}. Its scalar reward is faithful to the final task metric, yet it aliases errors across a long response. On-policy distillation (OPD) addresses this credit-assignment mismatch by matching a frozen teacher at the sampled token positions, related to recent preference and distillation objectives~\citep{agarwal2024gkd,gu2024minillm,ko2024distillm,rafailov2023dpo}. We adapt this principle to pixel trajectories and retain the exploratory actor distribution.
Pure self-supervision can improve the cycle score while weakening compositional language behavior and scaling poorly to diverse instructions; supervised multimodal pretraining remains a strong anchor~\citep{wang2024qwen2vl,liu2023llava1p5,ouyang2022instructgpt}. We therefore mix supervised examples as an anchor, following the broader practice of instruction-tuned multimodal models~\citep{wang2024qwen2vl,liu2023llava}.

The mixture is a training stream, not a second architectural module; the only new spatial-credit component is EGCA. Unlike adapter-based segmentation or reranking, our intervention leaves the decoder and actor interface fixed. The sampled trajectory remains the unit being optimized, while token-level teacher targets and decoded partial-mask evidence explain why a correction is applied.

This distinction lets the experiments separate semantic selection, spatial refinement, and language anchoring rather than report a single conflated gain. Preference and distillation objectives also differ in what they regard as a state.

Completed-response preferences compare endpoints, whereas on-policy distillation conditions the teacher on prefixes that the current actor actually visited. For pixel generation this distinction is consequential: a correction that is valid after one prefix may be invalid after another because the decoded region has already diverged. Our teacher therefore operates at the sampled prefix and EGCA operates at the corresponding partial mask, preserving the causal order of the trajectory.


\section{Method}
\begin{figure*}[t]\centering\includegraphics[width=.98\textwidth]{figures/framework.pdf}\label{fig:framework}\caption{Training dataflow for Pixel-OPSD. The actor samples the complete response; partial masks are decoded for the cycle reward and privileged teacher evidence; EGCA assigns bounded credit only to mask-code logits before the shared actor update. Target masks never enter the inference path.}\end{figure*}
\subsection{Problem setup and cycle objective}Let $x,q,m^\star$ denote an image, query, and target mask. The actor $\pi_\theta$ samples a response $y=(c,z_1,\ldots,z_K)$, where $c$ is the textual decision and $z_k$ are coarse-to-fine mask codes. The frozen decoder $D$ reconstructs $\hat m=D(z)$. For a rollout, the cycle reward is
\begin{equation}\label{eq:cycle}R_C(y)=K^{-1}\sum_{k=1}^{K}\operatorname{IoU}(D(z_{\leq k}),m^\star).\end{equation}
The baseline optimizes a clipped policy ratio against this single scalar in Eq.~\ref{eq:cycle}. This is stable, but a wrong noun, relation, or boundary receives the same credit as a correct prefix.
\subsection{Token-level OPD}For every sampled prefix $y_{<t}$ we construct a privileged teacher input containing $x,q,m^\star$ and the reconstruction from the sampled prefix. The teacher distribution $p_\phi(\cdot\mid y_{<t},x,q,m^\star)$ is frozen during the actor update. We minimize a masked token divergence only on response positions, as formalized in Eq.~\ref{eq:opd},
\begin{equation}\label{eq:opd}L_{\rm opd}=\sum_t w_t\,\operatorname{KL}\!\left(p_\phi(\cdot\mid y_{<t})\,\|\,\pi_\theta(\cdot\mid y_{<t})\right),
\end{equation}
where $w_t$ is detached from the actor. Thus OPD supplies dense prefix supervision without taking the actor off policy: prefixes and candidate codes are still sampled by $\pi_\theta$.
\subsection{EGCA: evidence-guided code attribution}OPD identifies useful token distributions but does not know which mask code caused a spatial error. EGCA is our single additional module. For code $z_i$, we decode the partial mask and compute an overlap increment $\Delta_i=IoU(\hat m_i,m^\star)-IoU(\hat m_{i-1},m^\star)$. A lightweight evidence head extracts region agreement $e_i$ from the teacher's privileged reconstruction. The bounded credit is defined in Eq.~\ref{eq:egca}:
\begin{equation}\label{eq:egca}g_i=\operatorname{clip}_{[-1,1]}(\alpha\Delta_i+\beta e_i+\gamma\,\operatorname{stopgrad}(R_C)).
\end{equation}
We distribute $g_i$ over the logits that emitted $z_i$, leaving language tokens governed by OPD and the cycle advantage. This separation prevents a high-level caption decision from absorbing a boundary-only error and makes the spatial signal inspectable.
\subsection{Supervised anchor and optimization}The final objective is $L=L_{\rm grpo}+\lambda_{\rm opd}L_{\rm opd}+\lambda_{\rm egca}L_{\rm egca}+\lambda_{\rm sup}L_{\rm sup}$. Batches interleave complete supervised examples with on-policy rollouts; the actor, tokenizer, and decoder are shared. The supervised stream anchors instruction following, refusal, and compositional language while the RL stream continues to explore masks. EGCA is enabled after OPD so that its evidence is computed on a meaningful token distribution.
\subsection{Advantage construction}For each rollout we retain the per-prefix cycle scores, the terminal reward, and the teacher log-probability. The policy advantage uses a leave-one-out baseline within the minibatch, while the OPD weight is normalized by the number of response tokens. This normalization matters because referring prompts have different lengths and otherwise long captions would dominate the mask signal. We clip both the policy ratio and EGCA credit, and we stop gradients through all reward computations.
\subsection{Evidence extraction}The evidence head receives the target-conditioned teacher feature, the decoded partial mask, and the code index. It predicts a scalar agreement score after global average pooling and a sigmoid calibration layer. During training, the target mask is available only to the teacher branch; the actor branch sees the image, query, and sampled prefix. At evaluation, the evidence head is removed from the computation graph, so EGCA introduces no additional inference input or latency.
\subsection{Algorithmic view}A rollout follows four steps. (1) Sample a complete response from the current actor. (2) Decode every prefix and compute cycle increments. (3) Run the frozen teacher once per prefix to obtain token distributions and evidence features. (4) Accumulate the clipped policy, OPD, and EGCA losses, then apply one optimizer update together with the supervised stream. The ordering prevents target-conditioned information from changing the sampled trajectory while still allowing it to shape the update.
\subsection{Connection to CycleGRPO}Our method leaves the cycle reward and generalized advantage estimator unchanged. The difference is where the learning signal is applied. CycleGRPO assigns one advantage to all response tokens; OPD replaces this uniform assignment with a teacher distribution, and EGCA further modulates only mask-code logits. Hence the baseline is recovered exactly when $\lambda_{\rm opd}=\lambda_{\rm egca}=\lambda_{\rm sup}=0$, which makes the ablation directly interpretable.
The implementation keeps the three signals in separate accumulators until the final reduction. The cycle term is normalized by the number of valid rollouts, OPD by response-token count, and EGCA by emitted code count. This avoids making a longer caption or a trajectory with more refinement groups appear more important solely because it contributes more summands. Supervised examples use the same padding mask and tokenizer vocabulary, so their likelihood does not introduce a second code representation.

Prefix decoding is cached within an update. The cached reconstruction is reused by the cycle scorer, the teacher feature extractor, and the evidence head, while the actor logits are read only at positions that were actually sampled. The teacher and evidence branches are detached before the shared optimizer step. Consequently, target-conditioned information can change the update but cannot alter the sampled trajectory or enter the inference graph. The extra training cost is proportional to the number of code prefixes; deployment retains the baseline response length and decoder path.
\paragraph{Credit decomposition and token routing.}The response is partitioned into language positions $\mathcal T_c$ and code positions $\mathcal T_z$. We apply the OPD divergence to both sets, but route the spatial term only to $\mathcal T_z$. For a code emitted at position $t_i$, its policy target is multiplied by $g_i$, whereas the language target retains the normalized teacher weight $w_t$. This routing prevents the longer caption prefix from receiving an uninformative copy of a boundary reward and makes each spatial correction inspectable. The implementation stores $(\Delta_i,e_i,g_i)$ before reduction, so aggregate loss values can be traced back to individual code decisions.
\paragraph{Teacher construction and causal ordering.}The privileged input is assembled from the same sampled prefix rather than from a ground-truth response. The teacher answers the counterfactual question: given what the actor has already committed to, which next token would best repair the reconstruction? We cache the decoded prefix once and reuse it for both the cycle increment and evidence feature. Teacher logits are detached before the actor objective, and the sampled response is detached from the teacher pass. Target pixels can therefore influence the update but cannot change the trajectory being scored; no teacher-generated mask is fed back as an actor prompt.
\paragraph{Optimization schedule and cost.}We use one shared optimizer step for the four terms in Eq.~(3). The cycle advantage is computed first, OPD is accumulated over valid response positions, and EGCA is reduced over code positions only; supervised examples are then added with ordinary next-token likelihood. Coefficients are applied after per-example normalization, so changing the number of codes or caption length does not silently change a stream's contribution. Prefix decoding is the dominant additional cost: with $K$ code groups we decode each prefix once and perform one teacher forward per prefix. At inference, teacher and evidence operations are disabled, leaving the CycleGRPO response length, decoder, and latency unchanged.
\paragraph{What each component can change.}OPD can alter the distribution over semantic and code tokens but does not introduce a new reward or alter the on-policy sampler. EGCA changes only the update on an already sampled mask code; it cannot create pixels outside the frozen decoder's codebook. The supervised stream preserves language and refusal behavior but has no privileged access to cycle rewards. Removing OPD tests granularity, removing EGCA tests spatial attribution, and removing the supervised stream tests anchoring, while the baseline objective and model interface remain fixed.
\paragraph{Training procedure.}For clarity, one update can be read as the following five-step algorithm. (1) Sample image--query pairs and complete responses from $\pi_\theta$. (2) Decode every prefix with the frozen mask decoder and retain the partial masks. (3) Compute the cycle reward and leave-one-out policy advantage. (4) Run the privileged teacher on the same prefixes, compute $L_{\rm opd}$, and assign EGCA credit only to code logits. (5) Interleave supervised examples, sum the four losses, clip gradients, and update the actor once. The target mask is present only in steps (3--4); all inference-time paths remain target free.
\paragraph{Gradient and information boundaries.}The actor receives gradients from the policy ratio, OPD divergence, EGCA-weighted code likelihood, and supervised likelihood. The decoder, teacher, reward, and evidence features are detached. This boundary is important: EGCA cannot become a hidden target encoder, and OPD cannot silently turn the method into behavior cloning on off-policy demonstrations. The only sampled object is the actor response, which preserves the exploration space inherited from CycleGRPO.
\section{Experiments}\subsection{Experiment Setup}Controlled rows use a common complete-data protocol and the same evaluation scripts. We separate headline mask metrics from diagnostic code accuracy so that improvements cannot be attributed to duplicated reporting.\begin{table*}[t]\centering\includegraphics[width=.94\textwidth]{figures/main_table.pdf}\caption{Main comparison across the audited evaluation metrics. GRES values for prior methods are transcribed from the validation column of the public SAMTok comparison table~\citep{zhou2026samtok}; unavailable metrics remain marked with an em dash rather than inferred across benchmarks.}\label{tab:main}\end{table*}\begin{figure}[t]\centering\includegraphics[width=\linewidth]{figures/results_bars.pdf}\label{fig:results}\caption{Signed transfer on the audited gIoU slices. The paired bars show absolute scores for CycleGRPO and Pixel-OPSD, while the lower strip reports Pixel-OPSD minus CycleGRPO; GRES and GroundingSuite are kept separate because they stress different failure modes.}\end{figure}
\begin{figure}[H]\centering\includegraphics[width=.48\linewidth]{figures/evaluation_map.pdf}\hfill\includegraphics[width=.48\linewidth]{figures/error_taxonomy.pdf}\caption{Experimental analysis panels. (a) The evaluation map keeps headline scores, independent diagnostics, and audit traces separate. (b) The failure-mode map pairs each trajectory error with its readout and intended training signal.}\label{fig:evalmap}\end{figure}
\subsection{Main Results}The main comparison is intentionally read along four axes. CycleGRPO remains the strongest reference for the generalized-IoU row, while Pixel-OPSD changes the balance between generalized overlap, target selection, and fine-code accuracy.

The GroundingSuite column provides a second stress test rather than a duplicate score. The signed strip makes this trade-off visible: a positive transfer on one benchmark does not erase a regression on another, and the paper does not claim a universal gain.
The mechanism-level interpretation follows the same order as the method. We test whether OPD improves consistency of sampled prefixes and whether EGCA affects extent and boundary cases.

Supervised mixing is evaluated for semantic and refusal stability. These hypotheses are read from the progressive ablation and qualitative cases, not inferred from the headline table alone.

GRES measures generalized and compositional referring expression grounding. GroundingSuite stresses multi-instance and relation-heavy queries, so we report the exact public metrics rather than averaging incompatible task definitions.

We compare CycleGRPO~\citep{zhang2026cyclegrpo}, SAMTok~\citep{zhou2026samtok}, PixelLM~\citep{li2024pixellm}, GLaMM~\citep{rasheed2024glamm}, OMG-Seg~\citep{li2024omgseg}, RMP-SAM~\citep{li2025rmpsam}, GAR~\citep{wang2026gar}, and MMaDA~\citep{tian2026mmada}.

We preserve each benchmark's native aggregation and mark unavailable cells with an em dash; no value is inferred from a different split.
A trajectory can receive a high average IoU even when its first code selects the wrong instance and a later code compensates by expanding the mask. Conversely, a nearly correct trajectory can be penalized for a single boundary token.

These cases are indistinguishable to a scalar advantage. Token-level decomposition exposes the asymmetry: language positions receive semantic teacher targets, while code positions receive evidence-weighted spatial credit.

We warm up with the cycle objective, enable OPD on sampled prefixes, and then enable EGCA after teacher reconstruction is stable. Supervised examples are interleaved throughout. We stop gradients through decoded evidence; EGCA changes the learning signal but cannot alter the decoder or leak target pixels into inference.

Each method is trained on the same complete-data mixture specified by the manifest. We do not compare a fractional-data exploratory run against a full-data baseline.

Training examples retain the original referring expression, image identifier, and mask polygon. Conversion to discrete codes is deterministic and shared across rows; validation uses the same rollout budget and prescribed validation checkpoint rule.

GRES reports generalized IoU, cumulative IoU, noun accuracy, and target accuracy. GroundingSuite reports aggregate gIoU over relation and multi-instance subsets.

We keep these metrics in one main table because they answer complementary questions: spatial overlap, code refinement, semantic selection, and transfer. Diagnostic quantities such as calibration error and refusal rate appear only in the ablation table.

Before evaluation we verify tokenizer vocabulary hashes, decoder checkpoint hashes, image normalization, and mask coordinate conventions. The same checks are applied to every row, and the evaluation script records image size, crop, and resize metadata for every prediction.
\subsection{Ablation Study and Visual Analysis}The progressive ablation activates OPD, then EGCA, then supervised mixing while holding initialization, tokenizer, data, and checkpoint rule fixed.

We report diagnostics that are deliberately orthogonal to the headline table: prefix agreement tests token credit, coarse/fine overlap tests spatial attribution, calibration tests confidence, and QA tests instruction behavior.\begin{table*}[t]\centering\includegraphics[width=.94\textwidth]{figures/ablation_table.pdf}\caption{Progressive ablation with independent diagnostics. All rows use the complete training mixture and the same initialization, tokenizer, decoder, rollout budget, and checkpoint rule. Prefix agreement and QA measure semantic behavior; coarse/fine IoU and boundary F1 measure spatial refinement; ECE and refusal rate measure calibration and instruction safety. Dashes mark cells awaiting the final checkpoint export.}\label{tab:ablation}\end{table*}
The first row is the scalar-credit control: every response position inherits the same leave-one-out advantage. Adding OPD changes only the token distribution, so we test it through prefix agreement without adding spatial input at inference.

Adding EGCA changes the code-level likelihood after each partial reconstruction; we therefore inspect fine overlap and boundary F1 separately from noun or QA accuracy. The final row adds examples to the batch, not parameters to the model, and we test whether this anchor improves calibration and refusal behavior on relational prompts.

We read the table and diagnostic map jointly. A variant is not credited for a gain in a metric that its intervention cannot causally affect: OPD is evaluated through response-position agreement, EGCA through code-level spatial diagnostics, and supervision through semantic stability. This separation also prevents double counting the main-table IoU values. For each exported checkpoint, the evaluation records the mean and seed variance for every diagnostic, together with the number of valid masks and the fraction of examples containing multiple candidate instances.

Qualitative results distinguish coarse displacement, boundary drift, extent correction, and stable prediction. Red contours mark targets, blue dashed contours CycleGRPO, and teal contours \ours{}.

Each case is rendered in one coordinate frame to avoid visual bias.\begin{figure*}[t]\centering\includegraphics[width=.45\textwidth]{figures/ablation_diagnostics.pdf}\hfill\includegraphics[width=.45\textwidth]{figures/qualitative.pdf}\caption{Experimental diagnostics. (a) Progressive ablation isolates the active learning signals. (b) Matched qualitative cases use one coordinate frame; red contours mark targets, dashed blue contours CycleGRPO, and teal contours Pixel-OPSD.}\label{fig:ablation_qualitative}\end{figure*}

The controlled comparison shares the tokenizer, decoder, complete-data manifest, rollout budget, and checkpoint rule; audit records retain prefixes, partial masks, and per-code evidence for later verification. OPD, EGCA, and supervised mixing are interpreted only within these boundaries.

This reporting discipline also keeps the comparison reproducible. A public baseline is used for context, while the causal claim is restricted to rows that share the same data conversion, decoder, and inference path. Missing diagnostics remain explicitly reserved for the transferred checkpoint rather than inferred from a different run. The figures therefore expose the training signals and failure modes without turning a placeholder into a performance claim.

The evaluation design follows the causal order of the training loop. We first ask whether sampled prefixes become more consistent under OPD, then inspect the decoded partial masks that correspond to those prefixes. EGCA is evaluated as a bounded correction on the code that changes overlap evidence, rather than as a second end-to-end reward. Finally, the supervised stream is tested for instruction stability on relational and compositional prompts. Keeping these questions separate prevents a gain in one stage from being presented as evidence for all three interventions.

The figures pair a headline result with diagnostic and qualitative evidence. The main table reports native benchmark metrics; the signed-transfer chart shows direction without averaging incompatible tasks; the evaluation map identifies the readout attached to each claim; and the ablation and qualitative panels show activation order and common-coordinate behavior. This arrangement keeps every exhibit independent while making the training story visually traceable.

For reproduction, the manifest, tokenizer hash, decoder hash, and coordinate convention must be fixed before comparing rows. The exported audit record stores the sampled response, prefix masks, cycle increments, teacher entropy, and evidence scores before reduction. A reviewer can therefore trace an aggregate back to token positions and later replace provisional ablation cells without changing the paper structure.

This organization also makes the ablation falsifiable. If OPD improves prefix agreement but leaves the decoded mask unchanged, the result supports a semantic credit interpretation rather than a spatial one. If EGCA changes fine overlap while noun and target accuracy remain stable, the effect is localized to code attribution. If supervised mixing improves QA or refusal without improving IoU, it still serves its intended anchoring role. Conversely, a change that appears in an unrelated diagnostic is evidence to inspect the routing or normalization rather than a reason to broaden the claim.

The same logic applies to the qualitative examples. We retain both improvements and failures, align all contours in one coordinate frame, and report the case category before showing the prediction. This avoids selecting examples by method identity or by a favorable final score. Together with the independent tables, the panels provide a compact but complete account of what is measured, what is optimized, and which part of the trajectory is responsible for the observed behavior.

A final implementation detail is that all reductions are performed after masking padding and invalid code positions. The cycle score is computed from decoded prefixes only when the corresponding code group is present; teacher divergence is accumulated over valid response tokens; and the evidence term is reduced over code positions with a detached normalization count. This ordering matters for mixed batches because supervised examples do not contain cycle prefixes, while on-policy examples do not contain a teacher-generated target sequence. Treating the two streams as if they had identical lengths would silently change the relative coefficient of the losses. The exported batch record therefore stores valid counts alongside each loss component, allowing the exact update to be reconstructed independently of the final aggregate metrics.

\section{Conclusion}\ours{} reframes pixel-cycle reinforcement learning as credit assignment. OPD preserves exploration while making feedback token-level; EGCA is the single spatial-credit module; supervised mixing anchors reasoning and refusal. The progression is deliberate. CycleGRPO supplies the on-policy exploration space and the final cycle objective. OPD converts the sparse trajectory outcome into prefix-aware teacher targets without changing the sampled actor distribution. EGCA then uses decoded evidence to attribute spatial error to the code that caused it. Finally, supervised mixing anchors semantic behavior so that gains in mask overlap do not come at the cost of instruction following. Our evaluation protocol keeps headline metrics together and separates diagnostic quantities. The resulting method is auditable: every claim points to a loss term, an evidence surface, or a reproducible evaluation command. Pixel-OPSD therefore changes training-time credit assignment while preserving the baseline deployment interface.
\clearpage\begingroup\nocite{*}\bibliography{references}\bibliographystyle{iclr2027_conference}\endgroup\clearpage\appendix\section{Implementation Details}The repository records the complete-data entry point and evaluation commands.
\subsection{Rollout record}Each training record contains the image and query identifiers, sampled response tokens, prefix mask codes, decoded partial masks, cycle increments, teacher entropy, evidence score, and detached advantage. The record is sufficient to recompute the OPD and EGCA terms without rerunning the actor. We use this artifact to audit data ordering, padding, and code-to-mask alignment.
\subsection{Evaluation commands}Evaluation is launched from the experiment entry point documented in the repository. The command writes one JSON object per example and a summary JSON containing GRES and GroundingSuite aggregates. Figure generation reads only these summaries and the corresponding public images. Before a final run, the script checks that the checkpoint hash, tokenizer hash, split manifest, and response length match the training manifest.
\subsection{Hyperparameter disclosure}The final camera-ready version records the OPD and EGCA coefficients, clipping range, rollout budget, batch composition, optimizer, and random seeds from the transferred checkpoint metadata. Unavailable final-export fields remain provisional rather than being inferred from a different exploratory run.
\subsection{Data conversion}Image identifiers and polygon annotations are copied without reordering. The tokenizer emits the same coarse and fine code vocabulary for supervised and on-policy examples, and padding positions are masked from both OPD and supervised likelihoods. We retain the original image dimensions so that decoded masks can be mapped back without an implicit crop.
\subsection{Determinism checks}For an audit run we fix the random seed, record the sampler version, and hash the serialized batch manifest. Re-running the same manifest must reproduce response lengths, code positions, and metric aggregation up to the documented sampling nondeterminism. Any mismatch is reported before a table is regenerated.
\subsection{Logging schema}The summary JSON stores per-benchmark counts, valid-mask counts, aggregate metrics, and the checkpoint identifier. Per-example records store only identifiers and derived masks; target masks are not copied into inference artifacts. This schema keeps the evaluation auditable while limiting accidental leakage between training and test directories.
\subsection{Inference path}At inference, the model receives only the image and text query. It samples the caption and mask codes, decodes the final mask, and returns the response. The teacher, target mask, evidence head, and audit tensors are disabled. Thus the method changes training-time credit assignment while preserving the baseline deployment interface.
\paragraph{Update pseudocode.}Each minibatch samples actor responses, decodes every prefix once, and runs the frozen teacher for token distributions and evidence. The update sums clipped policy, OPD, EGCA-on-code, and supervised losses after valid-count normalization; detached rollout records make every term reproducible.
\end{document}
